% ==================
% 2) THEORETICAL FRAMEWORK
% ==================
\section{Theoretical framework}\label{sec:theory}

The collateral universe is the tokenised funds of Section~\ref{sec:overview-channel}, whose only exit is issuer redemption at the published valuation over contractual windows. The exit is a process rather than a trade. The framework therefore models no venue depth, and reserves the cost of exiting in stress through the stressed SIMM risk weights, the concentration risk factor, and the clearing discount instead. Any incidental on-chain liquidity counts as zero, out of prudence. The theory reduces to three objects: a price process that jumps only at valuation publications, a worst-case loss over the contractual exit horizon, and a set of closed-form limits derived from them.

Per collateral $a$, the recipe is
\begin{equation}\label{eq:recipe}
H^{(a)} = \mathrm{IM}^{(a)} + \delta_{\mathrm{clear}}^{(a)} + \delta_{\mathrm{carry}}^{(a)} + \Lambda^{(a)},\qquad
\theta^{(a)} = \kappa^{(a)}\bigl(1 - H^{(a)}\bigr),\qquad
\ell_{\max}^{(a)} = \theta^{(a)} - \nu^{(a)},
\end{equation}
(the tier multiplier $\kappa^{(a)}\le 1$, a free parameter of the practitioner's tier structure, scales the range; $\kappa=1$ at the Prime tier, so $\theta=1-H$ there)
with caps $B_{\max}^{(a)}\le R_{\mathrm{cap}}^{(a)}$ and $S_{\max}^{(a)}\le M\,R_{\mathrm{cap}}^{(a)}$ (Section~\ref{sec:theory:capacity}). Here $\mathrm{IM}^{(a)}$ is the SIMM delta margin of Eq.~\eqref{eq:simm-core}, $\delta_{\mathrm{clear}}^{(a)}$ the clearing discount of Eq.~\eqref{eq:m-scaling}, $\delta_{\mathrm{carry}}^{(a)}$ the interest cost of Eq.~\eqref{eq:delta-carry}, $\Lambda^{(a)}$ the layer total of Eq.~\eqref{eq:layers}, and $\nu^{(a)}$ the per-listing liquidation buffer of Section~\ref{sec:policy-notch}. Every input is a published SIMM parameter, an attested redemption term (the issuer- or transfer-agent-attested notice, dealing-window, and settlement-lag figures of Section~\ref{sec:theory:setting}), an observed clearing cost, or a fixed policy constant. Nothing is fitted to market data that does not exist for this universe.

\subsection{Setting and notation}\label{sec:theory:setting}

Fix a collateral $a$. The official published valuation per token $N_k^{(a)}$ is published by the administrator at publication times $t_0 < t_1 < \cdots$ with publication cycle $\Delta^{(a)} = t_{k+1}-t_k$ set by the attested publication calendar: daily for a typical money-market fund (MMF), up to monthly for private credit and hedge funds. Between publications the protocol price accrues deterministically at the attested yield $y^{(a)}$, the portfolio yield net of fees, attested by the administrator alongside the valuation feed:
\begin{equation}\label{eq:nav-process}
P_t^{(a)} = \tilde N_t^{(a)} = N_k^{(a)}\bigl(1 + y^{(a)}(t-t_k)\bigr),\qquad t\in[t_k,t_{k+1}),
\end{equation}
and all price risk realises as a jump at the next publication,
\begin{equation}\label{eq:nav-jump}
J_{k+1}^{(a)} = \frac{N_{k+1}^{(a)}}{\tilde N^{(a)}_{t_{k+1}^-}} - 1 .
\end{equation}
The assets in scope have no continuous secondary-market price: intraday variance, order-flow impact, and oracle-versus-market deviation are not defined objects for this process. A published valuation does not by itself imply redeemability: gates, suspensions, and side-pockets can narrow the exit while valuations continue to be published. That state has no object in this price process. It is charged by the gating/suspension layer of Section~\ref{sec:theory:layers} and managed by the gate procedures of Section~\ref{sec:operations}.

Exiting a position means redeeming with the issuer. A forced exit cannot time the submission calendar, so the clearing window runs from the breach to cash: the worst-case wait $w$ to the next submission deadline, then the attested notice period $d_n$, dealing window $d_d$, and settlement lag $d_s$ (all in business days),
\begin{equation}\label{eq:mpor-eff}
\mathrm{LH}^{(a)} = \max\bigl\{1,\; w^{(a)} + d_n^{(a)} + d_d^{(a)} + d_s^{(a)}\bigr\} .
\end{equation}
$\mathrm{LH}^{(a)}$ is a contractual, attestable quantity rather than an estimate. For an asset dealing on fixed calendar dates the wait varies with the breach date, and the maximum is taken. Averages are never used.

Three sources fix the legs, in a fixed order. A written service-level agreement granted by the issuer or transfer agent in the attested redemption terms fixes them where one exists. It replaces the regulatory horizon rather than taking the larger of the two, so a shorter window is a reduction earned by evidence: a securitisation-type collateral with an attested redemption service-level agreement of 45 business days takes $\mathrm{LH}=45$ business days plus its worst-case wait. Where no such agreement exists but the issuer's own offering documents state a redemption or liquidity horizon, the framework reads $d_n$, $d_d$, and $d_s$ from those contractual statements, with the worst-case wait still added. Where neither exists, the regulatory liquidity horizon for the exposure under the Fundamental Review of the Trading Book (FRTB) \cite{bcbs_frtb_d457} applies: 40 business days for investment-grade credit, 60 for high yield, and 120 for securitisation. That row is floored at the contractual dealing window wherever that window is longer, and the worst-case wait is added to it as well, so the same collateral with neither takes $\mathrm{LH}=120$ business days plus its worst-case wait, the securitisation floor.

Where a class has no natural FRTB horizon, as with event-linked collateral, the practitioner sets a conservative fallback window as a model constant. A gating provision in the offering documents is charged separately in the gating layer rather than lengthening the window, so the contractual horizon and the gating risk are each charged exactly once. The clearing discount of Eq.~\eqref{eq:m-scaling} is built on this window.

\paragraph{Publication frequency and dealing frequency.} Many issuers publish valuations more often than they deal, a daily publication against monthly dealing being the common case, and the two frequencies are separate inputs. Position health is tested at every published valuation, so a breach can arrive on a publication when no redemption can be submitted. $\mathrm{LH}^{(a)}$ covers that case, so a breach mid-cycle is reserved on the same basis as one arriving on a dealing date. What follows a breach is lending-market behaviour, which the framework takes as given. Controlled operational records state when position health is evaluated and how a breach between dealing dates is treated.

\subsection{Worst-case valuation change over the clearing window}\label{sec:theory:stress}

The quantity the haircut must bound is the 99\% worst-case decline in redemption value of one unit of collateral between the breach and cash settlement, i.e.\ over the clearing window $\mathrm{LH}^{(a)}$. We do not estimate this from the token's own (short, appraisal-smoothed) valuation history. Instead we inherit from ISDA SIMM v2.6 \cite{isda_simm_v2_6} the prescribed measure of a 99\% move over ten business days, restricted to delta margin. The administrator attests a vector of sensitivities $s_k$ in the Common Risk Interchange Format. That file carries DV01 and CS01 ladders, the sensitivities to a one-basis-point move in interest rates and in credit spreads respectively, together with equity and FX deltas, per unit of published valuation. The model uses administrator-attested sensitivities where they are available and applies the conservative proxy rules of Section~\ref{sec:tax-simm-assignment} otherwise. Each sensitivity is weighted by its published risk weight and concentration risk factor. The five risk classes are interest rate, qualifying credit, non-qualifying credit, equity, and FX, where non-qualifying credit covers securitisations and other non-qualifying names. The weighted sensitivities aggregate within buckets under the published correlations $\rho$, across buckets under $\gamma$ with the bucket sums $S_b$ bounded to $[-K_b,K_b]$, and across the five risk classes $r\in\{\mathrm{IR},\mathrm{CreditQ},\mathrm{CreditNonQ},\mathrm{Equity},\mathrm{FX}\}$ through the inter-class matrix $\psi$. That roll-up runs inside each SIMM product class $p$, and the product-class margins then add:
\begin{equation}\label{eq:simm-core}
\mathrm{WS}_k = \mathrm{RW}_k\, s_k\, \mathrm{CR},\qquad
\mathrm{IM}_p^{(a)} = \sqrt{\textstyle\sum_{r\in p} K_r^2 + \sum_{r\in p}\sum_{r'\ne r}\psi_{rr'}\,K_r K_{r'}},\qquad
\mathrm{IM}^{(a)} = \sum_p \mathrm{IM}_p^{(a)},
\end{equation}
where $K_r$ is the class-level margin built from the $\mathrm{WS}_k$ and the product-class margins take no diversification between them (Section~\ref{sec:agg-haircut}). The full engine is specified in Section~\ref{sec:normalisation}. Per-type bucket assignment is fixed in Section~\ref{sec:taxonomy}. The scope is delta margin only, and Section~\ref{sec:tax-market} records why. Sensitivities are attested directly by the administrator, never derived statistically from the token's return history (no principal-component or factor-proxy construction), consistent with the no-fitting principle of Eq.~\eqref{eq:recipe}. Crypto-native collateral is out of scope entirely (Section~\ref{sec:overview}). An administrator who does not attest a sensitivity does not get zero: unattested rows fall to the Residual bucket at its conservative floor and the parameters freeze against loosening (Sections~\ref{sec:theory:layers} and~\ref{sec:normalisation}), so refusing to provide Common Risk Interchange Format sensitivities costs LTV, never risk. The model uses the latest valid inputs. Controlled records hold input snapshots and expiry dates. The SIMM parameter set itself is pinned to the published ISDA version.

SIMM is calibrated to a 99\% move over ten business days, and $\mathrm{IM}^{(a)}$ enters Eq.~\eqref{eq:recipe} verbatim at that horizon, with no rescaling factor applied inside the SIMM engine. The horizon the haircut must cover is the governing clearing window $\mathrm{LH}^{(a)}$, and Eq.~\eqref{eq:m-scaling} takes the 10-day move to it. The inherited standard uses ten business days as its reference horizon. The position is warehoused from the breach through the whole clearing window to cash settlement, $\mathrm{LH}^{(a)}$ business days in total (Eq.~\eqref{eq:mpor-eff}). Reconciling the move over that window to the charge at ten business days is an attribution convention, adopted to keep the SIMM core reconcilable against the published standard, and it makes no operational claim about what happens inside the window. The method also defines a stress equivalent under the Basel FRTB liquidity horizons \cite{bcbs_frtb_d457}. The controlled reconciliation record holds the per-listing comparison.

\paragraph{The clearing discount $\delta_{\mathrm{clear}}$.} The clearing discount $\delta_{\mathrm{clear}}^{(a)}$ is the difference between the price risk of the governing window and the charge at ten business days. It is a component of $H^{(a)}$, so the over-collateralisation that funds it sits above the liquidation threshold. The framework sets this reserve and does not prescribe the clearing mechanism: it is silent on how, and by whom, a position in liquidation is cleared, and on how the discount is enforced or paid. Eq.~\eqref{eq:m-scaling} sets the value:
\begin{equation}\label{eq:m-scaling}
\mathrm{IM}^{(a)} + \delta_{\mathrm{clear}}^{(a)} =
\begin{cases}
\mathrm{IM}^{(a)}\sqrt{\tfrac{\mathrm{LH}^{(a)}}{10}} & \text{day one (no qualified clearing history),}\\[4pt]
\max\!\left(\mathrm{IM}^{(a)}\sqrt{\tfrac{\mathrm{LH}^{(a)}}{10}},\ \delta_{\mathrm{obs}}^{(a)}\right) & \text{once } \ge N_{\mathrm{obs}} \text{ qualified observations exist,}
\end{cases}
\end{equation}
where the right-hand side is the 99\% move over the governing clearing window, and $\mathrm{IM}^{(a)}\sqrt{\mathrm{LH}^{(a)}/10}$ is the SIMM move over ten business days, rescaled to that window by the square root of time. The clearing discount is the residual, $\delta_{\mathrm{clear}}^{(a)}=\mathrm{IM}^{(a)}\bigl(\sqrt{\mathrm{LH}^{(a)}/10}-1\bigr)$ day one. That residual is positive for a window longer than ten business days and negative for a shorter one. A negative value is a margin credit, and it comes from contractually enforceable terms, attested or documented, that fix the window itself (Section~\ref{sec:theory:setting}). Estimated legs and FRTB fallback windows never earn one. Money-market collateral dealing daily is the case in point: its attested terms settle in days, and the framework charges the move over those days instead of over SIMM's ten business days. The credit never exceeds the SIMM charge, because $\mathrm{IM}^{(a)}+\delta_{\mathrm{clear}}^{(a)}=\mathrm{IM}^{(a)}\sqrt{\mathrm{LH}^{(a)}/10}\ge 0$ at every admissible window.

The rescaled move is a floor on the total price risk the haircut charges for, never a ceiling: where at least $N_{\mathrm{obs}}$ qualified clearing observations exist and their conservative quantile $\delta_{\mathrm{obs}}^{(a)}$ exceeds it, the observed clearing cost governs. $\delta_{\mathrm{obs}}^{(a)}$ is a level and not an increment, the realised clearing cost per unit of published valuation. An observation ageing out of the applicable period can lower $\delta_{\mathrm{obs}}^{(a)}$, and with it the total price-risk charge. A reduction below the rescaled floor requires a shorter attested $\mathrm{LH}^{(a)}$, which lowers the floor itself, or a documented amendment to the framework. The controlled validation record sets the qualification count, quantile, observation period, and diversity requirement.

Adverse evidence tightens the system: observed clearing costs above the prevailing reserve, or repeated failures to clear, force the escalation of Section~\ref{sec:operations}. That escalation can widen $\mathrm{LH}$, raise the relevant layer, cut caps, or pause the asset.

The square root of time caveat attaches here. The scaling is exact only for i.i.d.\ increments, and credit spreads are autocorrelated and jump-prone, so the rescaling is anti-conservative at long horizons for the credit types. The private-credit default add-on layer of Section~\ref{sec:theory:layers} is the explicit compensation there. The securitised and market-neutral fund tails are charged instead through their elevated credit and leverage risk weights, through the rating guardrail's automatic tightening on loss of investment grade, and through the leverage add-on, all of Section~\ref{sec:taxonomy}. The same law rescales the move downward for a window inside ten business days, where a clearing cost that does not shrink with the window would make the credit anti-conservative. The named layers compensate for that risk.

\paragraph{Reasons for the split.} The day-one identity of Eq.~\eqref{eq:m-scaling} is the same price-risk calculation, rescaled to the clearing window by the square root of time, so the split costs nothing numerically. It is split out for two reasons. Observed clearing evidence can inform $\delta_{\mathrm{clear}}$ and nothing else in the haircut, so keeping it distinct from the SIMM core keeps that evidence attributable (the qualified-observation rule of Eq.~\eqref{eq:m-scaling}). $\mathrm{IM}$ stays verbatim ISDA SIMM at ten business days, so the validation reconciliation against the published standard holds unchanged.

The same discipline governs missing data. An unavailable input is assigned a prescribed conservative floor in the spirit of FRTB's non-modellable risk factor (NMRF) treatment, and never zero, since zero-filling rewards exactly the assets that produce data voids (Section~\ref{sec:normalisation}). For an unattested $\mathrm{LH}$, the floor is the documented worst-case window. The practitioner sets a conservative model constant as the floor for each layer that requires one.

\paragraph{The interest cost $\delta_{\mathrm{carry}}$.} Debt accrues interest for as long as the position is held, so warehousing a position in liquidation through the clearing window has a certain cost that no other component of $H^{(a)}$ charges for. SIMM has no funding-cost risk class, and the clearing discount of Eq.~\eqref{eq:m-scaling} rescales price risk alone, so the interest cost is charged as its own component of the haircut:
\begin{equation}\label{eq:delta-carry}
\delta_{\mathrm{carry}}^{(a)} = \theta^{(a)}\, r^{B,\max}\, \frac{\mathrm{LH}^{(a)}}{252},
\end{equation}
where $r^{B,\max}$ is the ceiling of the deployed market's borrow-rate model, the annual percentage yield that model reaches at full utilisation, quoted on a 365-day calendar basis because debt in a decentralised lending market accrues on every calendar day. The deployed market configuration supplies the ceiling as a static input rather than a live utilisation reading. The worked example of Section~\ref{sec:normalisation} runs at an illustrative ceiling rate. A kink rate qualifies only where no higher rate is attainable. The window is contracted in business days and spans $\mathrm{LH}^{(a)}\times 365/252$ calendar days, the ratio of the calendar year to the business-day year, so the ceiling is charged over the year fraction $\mathrm{LH}^{(a)}/252$. That share is taken simple rather than compounded, which is the conservative reading of an annual percentage yield at every window inside a year. The rate applies to the debt outstanding when liquidation begins, $\theta^{(a)}$ per unit of published valuation, rather than to the whole collateral value. This is the only place the interest cost is charged: it never enters $\mathrm{IM}^{(a)}$, the liquidation buffer $\nu^{(a)}$, or the gating layer $\lambda_{\mathrm{gate}}^{(a)}$ of Section~\ref{sec:theory:layers}.

The threshold $\theta^{(a)}$ itself depends on $H^{(a)}$, so the two resolve together in closed form. Writing $c^{(a)} = r^{B,\max}\,\mathrm{LH}^{(a)}/252$ for the interest cost per unit of debt over the clearing window,
\begin{equation}\label{eq:carry-solve}
H^{(a)} = \frac{\mathrm{IM}^{(a)} + \delta_{\mathrm{clear}}^{(a)} + \Lambda^{(a)} + \kappa^{(a)} c^{(a)}}{1 + \kappa^{(a)} c^{(a)}},
\qquad
\theta^{(a)} = \kappa^{(a)}\bigl(1 - H^{(a)}\bigr),
\qquad
\delta_{\mathrm{carry}}^{(a)} = \theta^{(a)}\, c^{(a)} .
\end{equation}
The numerator holds every non-interest component of $H^{(a)}$: a type with a deterministic fee addend books it inside $\delta_{\mathrm{clear}}^{(a)}$, so the identity $\delta_{\mathrm{carry}}^{(a)} = H^{(a)} - H_0^{(a)}$, with $H_0^{(a)}$ the haircut before the interest cost, is exact.

Each period of interest accrual is reserved exactly once. The liquidation buffer $\nu^{(a)}$ charges for the borrow rate before liquidation begins, across one interval between published valuations, and it protects the borrower from being pushed through $\theta^{(a)}$ by routine drift (Eq.~\eqref{eq:notch-buffer}). The interest cost $\delta_{\mathrm{carry}}^{(a)}$ charges for the same rate after liquidation begins, across the clearing window, and it protects the deployment's lenders while the position is warehoused. The two intervals do not overlap and the two parties differ. Neither term enters $\lambda_{\mathrm{gate}}^{(a)}$, which charges for the chance that a clearing window does not pay and the cost of the extension that follows, never the certain interest cost over the contractual window.

\subsection{Additive layers for structural blind spots}\label{sec:theory:layers}

The sensitivity-based core models smooth co-movement of modelled risk factors. Everything it structurally cannot see enters as a named additive layer. The total layer charge is
\begin{equation}\label{eq:layers}
\Lambda^{(a)} = \sum_{j} \lambda_j^{(a)},
\end{equation}
a sum of fixed percentages per type. The practitioner sets the model constants used to determine these values (Appendix~\ref{sec:free-parameters}). Each $\lambda_j$ is tied to a specific structural blind spot:
\begin{enumerate}[label=(\roman*)]
\item \textbf{Gating/suspension ($\lambda_{\mathrm{gate}}$).} SIMM has no liquidity risk class: the narrowing of the exit (gates, side-pockets, suspensions) has no risk factor. This is the central tail risk of lending against redemption-routed collateral. Funds gate when redemption is most needed.
\item \textbf{Valuation staleness/smoothing.} Sensitivities are computed off administrator-published valuations that are smoothed relative to traded prices. Risk weights calibrated on traded histories partially compensate (a strength of inheritance), but the residual exposure between publications needs its own charge.
\item \textbf{Default add-on (private credit collateral).} SIMM margins spread delta rather than jump-to-default. FRTB pairs its sensitivity charge with a default risk charge for precisely this reason. A $PD\times LGD$ add-on in the FRTB-DRC style covers the gap.
\item \textbf{Oracle pipeline integrity.} The valuation oracle is a data pipeline with no SIMM risk factor. A weak configuration, such as a manual or single-signed feed, has its own charge by configuration tier, zero at the Robust tier.
\item \textbf{Prepayment (amortising consumer collateral).} Borrower prepayment optionality on amortising consumer collateral shifts cashflow timing in a way the delta margin alone does not charge for. SIMM's native compensator would be vega and curvature margin, but those require attested implied-volatility sensitivities, and no options trade on non-traded consumer collateral from which to attest them. The framework instead evaluates tranche CS01 at the slower of the disclosed conditional prepayment rate and a stressed-extension floor, and books the residual optionality as this named layer (the \P\,9 substitution of Section~\ref{sec:tax-market}; Section~\ref{sec:tax-simm-assignment}).
\item \textbf{Manager/operational.} Wrapper, transfer-agent, and administrator failure modes have no market sensitivity.
\item \textbf{Wrong-way.} Co-realisation of collateral stress with the stress of the borrower base (e.g.\ private-credit borrowers who are crypto market-makers) is not a SIMM object.
\item \textbf{Leverage add-on (leveraged fund collateral).} Fund-level leverage scales tail risk beyond what the fund-unit risk weight reserves. The add-on is charged per turn of leverage above policy.
\item \textbf{Event risk (event-linked collateral).} A collateral whose dominant loss is a jump with no SIMM risk factor at all, such as an insurance catastrophe event, takes a named event layer set from attestable model documents. Event probability grows roughly linearly in exposure time with undiminished severity, so neither the margin over ten business days nor the square-root-of-time extension can hold it, and the layer is the designated home.
\end{enumerate}
The decomposition $H = \mathrm{IM} + \delta_{\mathrm{clear}} + \delta_{\mathrm{carry}} + \Lambda$ keeps every charge attributable: a reviewer can challenge any single number without reopening the rest. The final haircut is thus D-SIMM. The risk-factor reference of the supporting documentation records, factor by factor, the metric behind each of these charges, its evidence source, and the parameter it drives.

\subsection{Policy limits}\label{sec:theory:limits}

The liquidation threshold and borrow limit follow from \eqref{eq:recipe}:
\begin{equation}\label{eq:policy-limits}
\theta^{(a)} = \kappa^{(a)}\bigl(1 - H^{(a)}\bigr),\qquad
\ell_{\max}^{(a)} = \theta^{(a)} - \nu^{(a)},\qquad \nu^{(a)}\in[0.04,\,0.08].
\end{equation}
The per-listing liquidation buffer $\nu^{(a)}$ is the \emph{borrower-side} gap between the borrow limit $\ell_{\max}^{(a)}$ and the liquidation threshold $\theta^{(a)}$. The 4--8pp range in \eqref{eq:policy-limits} is a fixed policy constant. Per-listing values within it are set at listing (Section~\ref{sec:policy-notch}).

\paragraph{Solvency by construction.} Liquidation begins when the loan-to-value ratio crosses $\theta^{(a)}$, so the over-collateralisation cushion protecting the system at the point of liquidation is $1-\theta^{(a)} \ge H^{(a)} = \mathrm{IM}^{(a)}+\delta_{\mathrm{clear}}^{(a)}+\delta_{\mathrm{carry}}^{(a)}+\Lambda^{(a)}$, with equality at the Prime tier ($\kappa=1$). Whatever discount compensates the party that holds the position to the redemption point is paid out of this cushion. The discount is bounded by $\mathrm{IM}^{(a)}+\delta_{\mathrm{clear}}^{(a)}+\delta_{\mathrm{carry}}^{(a)}$, and the cushion $1-\theta^{(a)}\ge H^{(a)}$ is available above the liquidation threshold. The layer total $\Lambda^{(a)}$ remains once the bound is exhausted, strictly so wherever the layer total is positive. Day one, before any qualified clearing observation governs, a window inside ten business days makes the price-risk charge $\mathrm{IM}^{(a)}+\delta_{\mathrm{clear}}^{(a)}$ smaller than the verbatim figure at ten business days. That is the intended reading of an asset reaching cash inside SIMM's ten-business-day horizon, and the interest cost $\delta_{\mathrm{carry}}^{(a)}$ is charged separately.

Placed below the liquidation threshold $\mathrm{LTV}=\theta^{(a)}$ instead, in the gap between the borrow limit and the liquidation threshold, the clearing discount could draw only on the buffer $\nu^{(a)}$ at the point of liquidation. That buffer is smaller than $\delta_{\mathrm{clear}}^{(a)}$ wherever the margin charge is large and the clearing window long, as for a bucket-11 fund unit at the published SIMM risk weight for that bucket (19 at v2.6) over a 45 to 95 business-day window. That design would be insolvent at the point of liquidation for exactly the assets that need the discount most.

Between two usable published valuations a position's debt accrues at the borrow rate while the collateral oracle accrues at the attested yield $y^{(a)}$, and at the next publication the collateral may jump by $J^{(a)}$. A borrower at exactly $\ell_{\max}$ must not be pushed through $\theta$ by that interest accrual plus one ordinary adverse valuation publication, which requires
\begin{equation}\label{eq:notch-buffer}
\nu^{(a)} \;\ge\; \ell_{\max}^{(a)}\,
\frac{r^{B,\max}\,\Delta t_A^{(a)}/252 + \bigl|q_p\bigl(J^{(a)}\bigr)\bigr|}
{1 - \bigl|q_p\bigl(J^{(a)}\bigr)\bigr|},
\end{equation}
where $r^{B,\max}$ is the ceiling borrow rate of the deployed interest-rate model. The inequality is exact. Debt grows by the factor $1+r^{B,\max}\Delta t_A^{(a)}/252$ over the interval while the published valuation the ratio is measured against falls by $\bigl|q_p(J^{(a)})\bigr|$, so the requirement takes the denominator $1-\bigl|q_p(J^{(a)})\bigr|$. Dropping that denominator understates the floor by exactly $\bigl|q_p(J^{(a)})\bigr|$ times the floor itself. At a $15\%$ floor it is $2.7$ points, against a buffer of 4 to 8 points. The interval $\Delta t_A^{(a)}$ is the longest gap between usable published valuations in business days, a usable valuation being one the deployment's valuation oracle can act on. It is the attested publication frequency, the publication cycle $\Delta^{(a)}$ of Section~\ref{sec:theory:setting} counted in business days, plus the attested publication lag, and the debt accrues across it on the 365-day calendar basis of Eq.~\eqref{eq:delta-carry}. The requirement takes no credit for the attested yield $y^{(a)}$: the published-valuation rebase and the accrual of interest on the debt run on different frequencies, so a yield offset can be too large or too small at any given time, and neither error is safe. The requirement is checked per listing in the controlled buffer record (Section~\ref{sec:policy-notch}), and the adopted constants must dominate it there.

The jump term $q_p(J^{(a)})$, the lower $p$-quantile of the jump between consecutive published valuations, is a \emph{prescribed} per-type stressed-jump floor. It is policy, never fitted: the framework bans estimating it from the token's own valuation history. The practitioner sets the stressed-jump floor and quantile level $p$ as model constants (Appendix~\ref{sec:free-parameters}). Eq.~\eqref{eq:notch-buffer} motivates the size of $\nu^{(a)}$ and is never evaluated on live data. The liquidation buffer is the static per-listing constant of Section~\ref{sec:policy-notch}. Types with infrequent valuation publication and large jumps between published valuations take the wide end of the range, and types with daily published valuations the narrow end. A type whose prescribed floor exceeds its adopted buffer launches under the new-borrow guardrail of Section~\ref{sec:policy-guardrails} rather than with a wider buffer. The liquidation buffer $\nu^{(a)}$ absorbs routine dynamics. The haircut $H^{(a)}$ is the reserve against stress. The two are not interchangeable, which is why \eqref{eq:policy-limits} subtracts them in sequence rather than blending them.


Eligibility and tiering sit upstream of \eqref{eq:policy-limits}: the practitioner's due-diligence scorecard screens listing through eligibility failures and assigns a tier multiplier $\kappa\le 1$ from the practitioner's own tier structure (illustratively $\{1.00,\,0.90,\,0.75\}$ for Prime/Standard/Restricted). A bridge or cross-chain dependence is not one of those failures. It places the asset outside the framework's scope, and the asset is not scored at all (Section~\ref{sec:overview-channel}). Lower tiers hold the larger over-collateralisation cushion $H_{\mathrm{eff}}=1-\theta\ge H$. The borrow limit follows as $\ell_{\max}=\theta-\nu$, with the liquidation buffer $\nu$ applied as an absolute deduction rather than scaled by $\kappa$. Section~\ref{sec:policy} states the public policy constraints. A fully worked composition of \eqref{eq:recipe} appears in Section~\ref{sec:normalisation}.

\subsection{Capacity and caps}\label{sec:theory:capacity}

The capacity reference, the single attested quantity from which both caps are set, is $R_{\mathrm{cap}}^{(a)}$, the issuer's attested redemption capacity in USD per dealing window. It is applied below and in Sections~\ref{sec:policy}--\ref{sec:efficacy} as the one-window capacity, the volume the issuer redeems within one dealing window. Where the dealing window differs from the basis on which the capacity is attested, the attested figure is window-normalised before it is used anywhere, because a capacity attested over a period longer than the dealing window overstates what one window can clear. The window basis is part of the definition of $R_{\mathrm{cap}}^{(a)}$ rather than a free parameter (Appendix~\ref{sec:free-parameters}). Both caps are closed-form functions of it:
\begin{equation}\label{eq:theory-caps}
B_{\max}^{(a)} \;\le\; R_{\mathrm{cap}}^{(a)},
\qquad
S_{\max}^{(a)} \;\le\; M\, R_{\mathrm{cap}}^{(a)},
\end{equation}
where $M$ is the supply-cap multiple: accepted supply is capped at $M$ dealing windows of attested capacity, a free parameter of the practitioner (Appendix~\ref{sec:free-parameters}). $B_{\max}^{(a)}$ caps borrowed \emph{debt} against collateral $a$ at the pool level and is denominated in units of borrowed debt, while $R_{\mathrm{cap}}^{(a)}$ and $S_{\max}^{(a)}$ are denominated in USD of collateral valued at the published valuation. The liquidator redeems collateral at the published valuation only until the debt is repaid, so the volume one dealing window must clear for the largest defaulting borrower is the \emph{debt} $B_{\max}^{(a)}$. The posted collateral above the repaid debt is redeemed later at the borrower's risk. Enforced at the pool level, the cap $B_{\max}^{(a)}\le R_{\mathrm{cap}}^{(a)}$ therefore guarantees that one window clears the pool's entire debt at the point of liquidation, and the interest accruing over the window is reserved separately by $\delta_{\mathrm{carry}}^{(a)}$. The supply cap bounds the aggregate portfolio in the same collateral units so that accepted supply clears within $M$ windows of attested capacity.

\paragraph{Control of the caps.} The deployment exposes caps per collateral rather than per entity, and the framework assumes no ability to identify a borrower: addresses are cheap, and nothing prevents one economic actor from borrowing through many of them. The enforced controls are therefore the per-collateral parameters a deployment may natively support: the supply cap $S_{\max}^{(a)}$, the borrow cap read at the pool level, and the ceiling $\ell_{\max}^{(a)}$. The deployed market configuration shows both active caps. Controlled records hold each recomputation from $R_{\mathrm{cap}}^{(a)}$ and the selected constant $M$. These support a guarantee that holds however debt is split across borrowers. Where the venue enforces the pool-level borrow cap at $B_{\max}^{(a)}\le R_{\mathrm{cap}}^{(a)}$, one dealing window clears the pool's entire debt at the point of liquidation. Where it enforces only the supply cap, total debt is at most $\ell_{\max}^{(a)}S_{\max}^{(a)}<M\,R_{\mathrm{cap}}^{(a)}$, so a full redemption completes within $M$ windows of dedicated capacity. The per-entity reading of $B_{\max}^{(a)}$, that the largest single borrower's debt clears in one window, is a monitored diagnostic rather than an assumption. Where the deployment cannot attribute debt to entities, the framework claims only the pool-level guarantee, which requires no identification of anyone.

\textbf{Remark.} The execution price does not vary with the quantity redeemed. Under redemption at the published valuation, the execution price of a redemption of size $q$ is the published valuation at settlement, constant in $q$ up to $R_{\mathrm{cap}}$, so $\partial P_{\mathrm{exec}}/\partial q = 0$. The mechanism that couples $\ell_{\max}$ to $S_{\max}$ in impact-based frameworks therefore has identically zero gain here: in those frameworks a larger portfolio forces a larger stress exit, which moves price further, which lowers the viable LTV, and that price move is absent under redemption at the published valuation. The remark is independent of the caps: $\partial P_{\mathrm{exec}}/\partial q = 0$ is a property of the redemption process alone, so the one-directional computation remains valid with per-entity concentration uncontrolled, as it is here. Caps and $\mathrm{CR}$ bound the time to exit and haircut steepness, never the exit price, so weakening them lengthens the feasibility horizon of Section~\ref{sec:feasibility} ($\lceil B_{(1)}/R_{\mathrm{cap}}\rceil$ windows to clear the largest debt instead of one) without reintroducing a circular dependency. The bound $\mathrm{CR}\le\sqrt{M}$ rests on the aggregate supply cap $S_{\max}$, a native parameter of the deployment unaffected by how per-entity concentration is controlled.

At each recomputation, the attested inputs of Eq.~\eqref{eq:recipe} give $\mathrm{IM}^{(a)}$, $\delta_{\mathrm{clear}}^{(a)}$, $\Lambda^{(a)}$, and $\delta_{\mathrm{carry}}^{(a)}$, then $H^{(a)}$, $\theta^{(a)}$, and $\ell_{\max}^{(a)}$. The attested capacity $R_{\mathrm{cap}}^{(a)}$, with the policy constant $M$, gives the borrow cap and the supply cap of Eq.~\eqref{eq:theory-caps}. No output is an input to its own calculation. $\mathrm{CR}$ is evaluated at the observed exposure, an input read from the current state rather than solved jointly with the limits. The re-keyed threshold $T=R_{\mathrm{cap}}^{(a)}$ and the bound $\mathrm{CR}\le\sqrt{M}$ are derived in Section~\ref{sec:normalisation}.

The borrow rate does not reopen a loop. The ceiling rate $r^{B,\max}$ is a static constant of the deployment, and the interest cost it reserves resolves in closed form at Eq.~\eqref{eq:carry-solve}. The rate a borrower actually pays appears only in the motivation for the buffer \eqref{eq:notch-buffer}, so a utilisation-dependent borrow rate feeds no live quantity dependence into the computation.
